\providecommand{\Kpre}{\mathsf{K}_{\mathrm{pre}}}
\providecommand{\sC}{\mathsf{C}}
\providecommand{\sK}{\mathsf{K}}
\providecommand{\para}[1]{\par\noindent\textbf{#1}\quad}
\providecommand{\todo}[1]{\textbf{[TODO: #1]}}

\subsection{Setup Details}
\label{app:eval-setup}

The dog-versus-cat ResNet-18 stands in for the unreleased classifier of Anchuri et
al.~\cite{anchuri2026verifiable}. Int8 quantisation changes the accuracy of the image models by at
most \CnnAccuracyChange{} points. The language models' costs use their real checkpoints (Section~\ref{sec:eval-setup}), as do the attack on
OPT-6.7B (Section~\ref{sec:break}) and model quality (Section~\ref{sec:quality}). Proofs are reproducible across
hardware: on \HwDigestCells{} configurations of the benchmark (image models and decoder blocks, modes~$\sC$ and~$\Kpre$,
v0 and V1, Fiat--Shamir), the prover on \HwDigestDevices{} produced byte-identical proofs and commitments.

The language models use 11-bit attention probabilities, as the certified models do. The scale is a
public constant of the attention, which the verifier recomputes: it changes no shape, only values, and with them the
compressed claims and, in mode~$\sC$, the columns that the challenges open. On two blocks of Llama-2-7B at
2{,}048 tokens (random weights, since a checkpoint's recipe fixes the scale), in modes~$\Kpre$ and~$\sC$ on the same queries with the CPU verifier on 8 threads, 11-bit
probabilities change the proof's size by at most \ProbBitsBytes{}\% and the prover's and verifier's medians by at
most \ProbBitsTime{}\% against 8-bit ones.

\subsection{Prover and Verifier Improvements}
\label{app:eval-improve}

\begin{table}[t]
  \caption{Times before and after each implementation improvement---each speeds up its phase and changes no claim, proof or verdict. One \PlatformGpu{}; the CPU row on 8 threads of an \PlatformCpu{}; $b/L$: $b$ of the model's $L$ blocks; $T{=}2{,}048$ tokens except for the int8 opening}
  \label{tab:eval-improve}
  \centering\footnotesize
  \setlength{\tabcolsep}{4pt}
  \begin{tabular}{@{}lrrr@{}}
    \toprule
    Build ($b/L$ blocks) & before & after & gain \\
    \midrule
    \pvitable{tab_eval_improve}
    \bottomrule
  \end{tabular}
\end{table}

Each row is measured on the \PlatformGpu{} with 8 threads of the \PlatformCpu{}, against the code path it
replaces on the same queries (Table~\ref{tab:eval-improve}); a row measures one improvement, except the GPU
verifier's row (its claim decoder and fused attention together) and the CPU verifier's row (all its native
kernels, int32 claims and the allocator setting together). No improvement changes a claim, a proof or a
verdict.

\para{Prover.} With float64 products, opening the committed columns takes \ImproveOpenShare{}\% of the prover's time
on Llama-2-7B blocks at 64 tokens (\ImproveOpenFloat{} of \ImproveProveFloat{}\,s). Splitting each field element into four signed bytes makes the
opening and the fold exact int8 GEMMs with bit-identical results, and the prover \ImproveOpenRange{}$\times$
faster on Llama-2 blocks. Int8 tensor cores also halve the forward pass's weight products on
Llama-sized layers, and the prover reuses the verifier's fused attention kernel. The claim encoder
was limited by a single-threaded copy of its output, so assembling the output on several threads
gives \ImproveEncodeRange{}$\times$. A fused Triton encoder that writes the same bytes is slower
(\ImproveFusedEncoder{} against \ImproveThreadedEncoder{}\,s on OPT-1.3B), so we do not use it. The prover
also frees each activation once its last consumer has run and keeps the int32 claims on the GPU until they
are encoded, which bounds its memory without changing its output.

\para{Verifier.} The GPU verifier path that these kernels replace spends most of its time decoding
the claims on the host and passing about a dozen times over materialised $T\times T$ attention
scores. Our GPU decoder validates every offset on the host before launching a kernel and rejects
exactly the inputs that the host decoder rejects (tested on every truncation and on 400 random
corruptions). Our fused
attention kernel computes the same integers without a $T\times T$ buffer; its fp32 sums are exact
because every partial sum is an integer below $2^{24}$. Together they make the GPU verifier
\ImproveGpuvRange{}$\times$ faster at 2{,}048 tokens. Afterwards, copies take \ProfileCopies{}\% of the CUDA time of an
OPT-1.3B query (12 blocks) and int8 GEMMs \ProfileIntMM{}\%: data movement dominates, including the upload of
the proof, which V1 makes smaller. On a CPU, native exact kernels for every phase (the recomputed operations, the AVX-512 attention, the
field products and the claim decoder), claims held as int32 instead of int64, and glibc's allocator set to
keep freed memory in the process together gain \ImproveCpuv{}$\times$ at 2{,}048 tokens.

\para{Bounded memory.} Two modes keep large cells within the GPU's memory without changing a byte or a
verdict, and the run uses them only where the default path does not fit: the GPU verifier of the Llama-2
cells at 2{,}048 tokens decodes the claims one stream at a time instead of all at once, and V1's
Llama-2-7B cells at 2{,}048 tokens keep the GKR's layers in bounded memory. Both run inside the timed region,
so the times of those cells include them.

\subsection{Memory and Untrusted Commitments}
\label{app:eval-untrusted}

\begin{table}[t]
  \caption{Memory of one Llama-2-7B query, in GB---the client needs no GPU, but its memory grows with the prompt. Client: peak resident set of the CPU verifier alone (8 threads of an \PlatformCpu{}) and after its imports; prover GPU: peak allocated memory of one \PlatformGpu{}; process: peak resident set of prover and verifier together}
  \label{tab:memory}
  \centering\footnotesize
  \setlength{\tabcolsep}{4pt}
  \begin{tabular}{@{}lrrrrr@{}}
    \toprule
    Mode & $T$ & client & after imports & prover GPU & process \\
    \midrule
    \pvitable{tab_memory}
    \bottomrule
  \end{tabular}
\end{table}

The prover and the verifier of a timed cell share one process, whose peak resident set is theirs
together (Table~\ref{tab:memory}, process). The client's own memory is therefore measured in a
separate process that holds only what the client holds (the public graph, its key and the
embedding tables it looks up in $\Kpre$ or the roots in $\sC$, and one proof at a time) and checks
every timed proof again from the proof's bytes; its verdicts equal those of the measuring process.
For Llama-2-7B in $\Kpre$ it peaks at \ClientPeakShortKpre{}\,GB at 64 tokens and at
\ClientPeakLongKpre{}\,GB at 2{,}048 tokens (\ClientPeakShortC{} and \ClientPeakLongC{}\,GB in $\sC$), of
which at most \ClientBaseline{}\,GB is the runtime after its imports (Python, PyTorch and our
library). Re-execution holds the \WeightsLlama{} of int8 weights.

Without the honest-commitment assumption (Section~\ref{sec:setup-proof}), Llama-2-7B needs a one-time
setup proof of \SetupGBLlamaA{}\,GB with parameters amortised over $a=100$ queries
(\SetupGBLlamaB{}\,GB for $a=1{,}000$). It raises the opened columns per interactive query from
\SetupHonestMBLlama{} to \SetupUntrustedMBLlamaA{}\,MB (\SetupUntrustedMBLlamaB{}\,MB), an increase of at most
\SetupQueryShareLlamaA{}\% of the proof at 2{,}048 tokens (Table~\ref{tab:setup}).

\subsection{Real Weights}
\label{app:eval-real}

Real weights cost about as much as random ones. The real graphs are the certified integer OPT
models of Section~\ref{sec:quality} (calibrated on four 2{,}048-token windows of the train split).
Each cell has a twin with random weights of the same shapes, and every cell runs \RealQueries{}
timed queries after a warm-up (three in OPT-6.7B's V1 cells), with the prover on the
\PlatformGpu{} and the native CPU verifier on 8 threads of the \PlatformCpu{}
(Table~\ref{tab:eval-real}). Every honest query was accepted, and every tampered one (a changed
claim, output logit or last hidden state) rejected.

\begin{table}[t]
  \caption{Real against random weights on OPT-1.3B---about the same time, and a smaller v0 proof with real weights. Each cell real / random; prover on one \PlatformGpu{}, verifier on 8 threads of an \PlatformCpu{}}
  \label{tab:eval-real}
  \centering\footnotesize
  \setlength{\tabcolsep}{2.5pt}
  \begin{tabular}{@{}lrrrrr@{}}
    \toprule
    Mode & $T$ & v0 (MB) & V1 (MB) & prove (s) & verify (s) \\
    \midrule
    \pvitable{tab_eval_real}
    \bottomrule
  \end{tabular}
\end{table}

The prover's time is within \RealProveWithin{}\% of the random twin's and the verifier's within \RealVerifyWithin{}\%,
while the timed queries of one cell spread by up to \RealQuerySpread{}\% of their median. v0's proof is
\RealProofSmaller{}\% smaller with real weights, because a calibrated model's claims are more concentrated than
a random one's, while V1's values take about the same space; V1's gain is therefore smaller (at 64
tokens \VOneRealOptMShortC{}$\times$ in $\sC$ and \VOneRealOptMShortK{}$\times$ in $\Kpre$, against \VOneRandomOptMShortC{}$\times$ and \VOneRandomOptMShortK{}$\times$).
In mode $\sC$ at 64 tokens, real weights cause \VOneRealOptMShortExceptions{} V1 exceptions (random: \VOneRandomOptMShortExceptions{}), which cost
\VOneRealOptMShortExceptionsMB{}\,MB.

The full real OPT-6.7B (32 blocks) runs end to end in $\Kpre$ at 2{,}048 tokens.
\RealOptLAccepted{} of \RealOptLQueries{} queries were accepted and \RealOptLTampersRejected{} of
\RealOptLTampers{} tampers rejected. The proof is \RealOptLProof{}\,GB (with embedding lookups); it
takes \RealOptLProve{}\,s to prove and \RealOptLVerify{}\,s to verify on 8 threads of the
\PlatformCpu{} (decoding \RealOptLDecode{}\,s, recomputation \RealOptLRecompute{}\,s, products
\RealOptLProducts{}\,s). On the V1 run's own queries, V1's proof is \VOneRealOptLVOneGB{}\,GB against
\VOneRealOptLVZeroGB{}\,GB for v0 (\VOneRatioRealOptLExact{}$\times$), with \VOneRealOptLExceptions{} exceptions
(\VOneRealOptLExceptionsMB{}\,MB); in mode $\sC$ V1 gains \VOneRatioRealOptLCExact{}$\times$ (random
twin: \RndOptLVOneRatioC{}$\times$). The random-weight twin takes \RndOptLProve{}\,s to prove and
\RndOptLVerify{}\,s to verify, with a \RndOptLProof{}\,GB v0 proof that V1 makes
\RndOptLVOneRatio{}$\times$ smaller. Real weights are thus within \RealOptLTimeWithin{}\% in time, and
their proof is \RealOptLVZeroSmaller{}\% smaller with v0 and \RealOptLVOneLarger{}\% larger with V1.

\subsection{Long Prompts}
\label{app:eval-long}

Beyond 2{,}048 tokens the full model is extrapolated from smaller builds. We measure builds of
Llama-2-7B's shapes with 2, 3 and 4 blocks (random weights) at 2{,}048 to 16{,}384 tokens ($\Kpre$, the
native CPU verifier on 8 threads of the \PlatformCpu{}, every query accepted) and exact
re-execution on the same threads (with glibc's default allocator, as in
Appendix~\ref{app:eval-cpu}), and extrapolate both to 32 blocks by a least-squares line through
the three builds (Table~\ref{tab:eval-long}). A block's verifier time follows $aT+bT^2$ within
\LongFitResidual{}\% from 2{,}048 to 8{,}192 tokens. The proof grows linearly, by
\LongBytesPerBlock{}\,GB per block at 8{,}192 tokens. 32{,}768 tokens are beyond the range of our
kernels.

The full model with random weights, as the fit (the long tier's cells), measured at 2{,}048 and 8{,}192
tokens,\footnote{At 8{,}192 tokens the full model's claims do not fit in the GPU's memory beside the model, so the prover
moves \LongCheckStaged{}\,GB of them to host memory as they are made. This changes only the prover's time, which this
check does not use; the verifier, the re-execution and the proof are as they would be without it.} checks
the extrapolation (Table~\ref{tab:eval-long-check}): the fit is within \LongCheckVerify{}\% of the measured verifier,
\LongCheckReexec{}\% of the measured exact re-execution and \LongCheckProof{}\% of the measured proof. A line through the 2- and
3-block builds alone is \LongCheckTwoThree{}\% off the verifier at 8{,}192 tokens. On the \PlatformGpu{}, the same fit through 2 and
3 blocks is within \LongCheckGpu{}\% of the full model's streaming verifier at 2{,}048 tokens.

Extrapolated from 2 and 3 blocks, the streaming GPU verifier needs about \LongGpuVerifyLong{}\,s for the full model at 16{,}384
tokens, and exact int8 re-execution on that GPU \LongGpuRatioLong{}$\times$ as long (Table~\ref{tab:eval-long-gpu}). Its fused attention keeps
no $T\times T$ score tensor: with 3 blocks at 16{,}384 tokens, the process that proves and verifies peaks at
\LongGpuPeakLong{}\,GB of GPU memory.

V1 keeps most of its gain at long prompts (Table~\ref{tab:eval-v1-long}). Whole V1 queries on 2 blocks at
2{,}048 to 8{,}192 tokens are accepted and their tampers rejected. A codec profile of the full model
(one forward pass, every claim encoded as v0 and as V1 send it) gives V1's claims ratio at each
length: \VOneClaimsRatioTwoK{}$\times$ at 2{,}048 tokens,
\VOneClaimsRatioEightK{}$\times$ at 8{,}192 and \VOneClaimsRatioSixteenK{}$\times$ at 16{,}384 (the profile's change with
length applied to the protocol's measured ratio at 2{,}048 tokens). With the ratio of each length, V1's proof of the full Llama-2-7B in mode $\sC$
reaches the size of its int8 weights at about \BreakevenTokensVOne{} tokens, interpolated between lengths measured up to
\BreakevenTokensVOneMeasuredTo{}; the 2{,}048-token ratio at every length would give about \BreakevenTokensVOneConstant{}.

\begin{table}[t]
  \caption{CPU verifier against exact re-execution at long prompts---the verifier stays faster, while the proof grows linearly. Full Llama-2-7B extrapolated from 2, 3 and 4 random-weight blocks; 8 threads of an \PlatformCpu{}; the fastest exact re-execution, glibc's default allocator}
  \label{tab:eval-long}
  \centering\footnotesize
  \setlength{\tabcolsep}{3pt}
  \begin{tabular}{@{}rrrrr@{}}
    \toprule
    $T$ & verifier (s) & re-execution (s) & ratio & proof (GB) \\
    \midrule
    \pvitable{tab_eval_long}
    \bottomrule
  \end{tabular}
\end{table}

\begin{table}[t]
  \caption{Measured against extrapolated values---the extrapolation is close to the full Llama-2-7B. Each cell measured / extrapolated; CPU: 8 threads of an \PlatformCpu{}, builds of 2, 3 and 4 blocks; GPU: the \PlatformGpu{}, builds of 2 and 3 blocks}
  \label{tab:eval-long-check}
  \centering\footnotesize
  \setlength{\tabcolsep}{3pt}
  \begin{tabular}{@{}lrrrr@{}}
    \toprule
    & $T$ & verifier (s) & re-execution (s) & proof (GB) \\
    \midrule
    \pvitable{tab_eval_long_check}
    \bottomrule
  \end{tabular}
\end{table}

\begin{table}[t]
  \caption{GPU verifier against GPU re-execution at long prompts---exact re-execution with the weights stays faster at every length. Full Llama-2-7B on the \PlatformGpu{}, extrapolated from 2 and 3 random-weight blocks; re-execution exact with int8 GEMMs, and fp16; memory: peak of the 3-block process that proves and verifies}
  \label{tab:eval-long-gpu}
  \centering\footnotesize
  \setlength{\tabcolsep}{3pt}
  \begin{tabular}{@{}rrrrrr@{}}
    \toprule
    $T$ & verifier (s) & int8 (s) & fp16 (s) & int8 / verifier & peak (GB) \\
    \midrule
    \pvitable{tab_eval_long_gpu}
    \bottomrule
  \end{tabular}
\end{table}

\begin{table}[t]
  \caption{v0 against V1 at long prompts---V1 keeps most of its gain. 2 Llama-2-7B blocks (whole queries, $\Kpre$, random weights); claims ratio: v0 over V1 for the full model, from its codec profile}
  \label{tab:eval-v1-long}
  \centering\footnotesize
  \setlength{\tabcolsep}{3pt}
  \begin{tabular}{@{}rrrrrrr@{}}
    \toprule
    & \multicolumn{5}{c}{2 blocks} & full model \\
    \cmidrule(lr){2-6}
    $T$ & v0 (MB) & V1 (MB) & ratio & prove (s) & verify (s) & claims ratio \\
    \midrule
    \pvitable{tab_eval_v1_long}
    \bottomrule
  \end{tabular}
\end{table}

\subsection{Catching the Edit, and the Setup Proof}
\label{app:eval-attack}

Full soundness also sends fewer bytes than making the path protocol catch the edit of
Section~\ref{sec:break}. Catching one edited neuron with probability $1-2^{-40}$ takes at least
\PathsLlama{} paths of the path protocol for Llama-2-7B, about \OpenAllLlama{}\,GB of openings at 64
tokens, where our proof is \ProofLlamaShortC{} at $\lambda = 128$ (Table~\ref{tab:eval-llm}). Without
the setup proof (Section~\ref{sec:setup-proof}), a commitment that mixes two models passes about
$2^{-t}$ of the queries, where $t$ is the number of opened columns; with it, such a commitment is
rejected at setup, as the setup proof's tests check. In our measurements, all
\OutcomeSetupProofs{} honest setup proofs were accepted. Besides the queries counted in
Section~\ref{sec:eval-setup}, the real-weight cells, their random-weight twins and the generation
cells sent \OutcomeTampers{} tampered queries (\OutcomeTampersVOne{} of them V1 queries and
\OutcomeTampersGen{} generated responses), and all were rejected.

\subsection{V1: Crossover per Cell}
\label{app:eval-crossover}

V1 pays off end to end only on slow links. Its prover takes about \VOneProverRatioAbout{}$\times$ as
long as v0's at 2{,}048 tokens, but still seconds where the zkSNARK provers of
Table~\ref{tab:eval-compare} take minutes, so V1 suits clients whose bandwidth matters more than
the prover's time. If proving, transfer and verification run one after
another, V1 is faster end to end on links slower than $8\,\Delta B/\Delta t$ bits per second, where
$\Delta B$ is the number of bytes it saves and $\Delta t$ the prover and verifier time it adds. On
the timed cells of Table~\ref{tab:eval-v1} (MB $=10^6$ bytes) this gives \CrossoverGptMediumC{}\,Gb/s (GPT-2, 512
tokens, $\sC$), \CrossoverGptMediumK{}\,Gb/s (GPT-2, 512 tokens, $\Kpre$), \CrossoverOptSLongC{}\,Gb/s (OPT-125M, 2{,}048 tokens,
$\sC$), \CrossoverOptSLongK{}\,Gb/s (OPT-125M, 2{,}048 tokens, $\Kpre$) and \CrossoverLlamaLongK{}\,Gb/s (Llama-2-7B,
\CrossoverLlamaLongKBuild{}, 2{,}048 tokens, $\Kpre$)\TodoCrossover{}.

\subsection{Re-Execution: Further Remarks}
\label{app:eval-reexec}

\begin{table}[h]
  \caption{Our verifier against re-execution and inference---on a CPU it is at least as fast as exact re-execution and faster than fp32 inference. Seconds, full models. GPU: one \PlatformGpu{}, int8 tensor cores and the fp16 decoder; CPU: Llama-2-7B on 8 threads of an \PlatformCpu{}; llama.cpp~\cite{llamacpp} (Q8\_0, bf16) computes a different model of the same shapes}
  \label{tab:eval-reexec}
  \centering\footnotesize
  \setlength{\tabcolsep}{4pt}
  \begin{tabular}{@{}lrrr@{}}
    \toprule
    & $T{=}64$ & 512 & 2{,}048 \\
    \midrule
    \pvitable{tab_eval_reexec}
    \bottomrule
  \end{tabular}
\end{table}

A CPU client that could keep the weights beats exact (integer) re-execution at every measured length
by keeping a $\Kpre$ key instead (Table~\ref{tab:eval-reexec}; \CpuRatioMedium{}$\times$ at 512 tokens). The llama.cpp rows run release
\LlamaCppRelease{}, built for the \PlatformCpu{}, on the same 8 threads, with Llama-2-7B's weights converted to
llama.cpp's own formats; each row is the fastest of its two attention kernels and of micro-batches of 512 and
2{,}048 tokens (at most the prompt), and, unlike our pruned graph, llama.cpp computes the last block's attention at
every position. Its faster format at 2{,}048 tokens, \CpuLlamaCppName{}, takes \CpuRatioLlamaCpp{}$\times$ our
$\Kpre$ verifier's time. Our streaming GPU verifier checks a 2{,}048-token OPT-1.3B query in $\Kpre$ in
\OptMGpuVerify{}\,s on the \PlatformGpu{}, against the 0.90\,s that zkLLM reports on an A100.

\subsection{The CPU Verifier: Cores, Kernels and Model Size}
\label{app:eval-cpu}

The ratio of Section~\ref{sec:eval-reexec} is measured on 8 cores of a server CPU.
Table~\ref{tab:eval-cpu} measures it for Llama-2-7B at 2{,}048 tokens on 4 and 2 cores of one NUMA
node, the verifier and re-execution each on as many threads, on 8 cores without our AVX-512
attention kernel, and for Llama-2-13B. Every re-execution in this table uses glibc's default
allocator, so that its rows compare like with like, whereas Section~\ref{sec:eval-reexec} takes the
fastest re-execution under either allocator setting; the 8-core ratio here
(\CpuRatioThrEight{}$\times$) therefore differs from the \CpuRatio{}$\times$ of that section. Our $\Kpre$ verifier is \CpuRatioThrFour{}$\times$ faster
than exact re-execution on 4 cores and \CpuRatioThrTwo{}$\times$ on 2, against
\CpuRatioThrEight{}$\times$ on 8. Without the AVX-512 attention, in the verifier and
in the integer forward pass alike, the verifier takes \CpuNoAvxSlowdown{}$\times$ as long and the ratio is
\CpuRatioNoAvx{}$\times$; the library is still built for the host CPU, so this isolates the attention kernel and
emulates no CPU without AVX-512. For Llama-2-13B, the verifier takes \CpuVerifyThirteen{}\,s against
\CpuReexecThirteen{}\,s for exact re-execution (\CpuRatioThirteen{}$\times$).

\begin{table}[t]
  \caption{CPU verifier against exact re-execution on fewer cores, without the AVX-512 attention kernel, and on Llama-2-13B---our verifier stays faster in every setting, most of all on few cores. Seconds at 2{,}048 tokens ($\Kpre$; the fastest exact re-execution, glibc's default allocator); cores of one NUMA node of an \PlatformCpu{}, one thread per core}
  \label{tab:eval-cpu}
  \centering\footnotesize
  \setlength{\tabcolsep}{4pt}
  \begin{tabular}{@{}llrrr@{}}
    \toprule
    Model & threads & verifier & re-execution & ratio \\
    \midrule
    \pvitable{tab_eval_cpu}
    \bottomrule
  \end{tabular}
\end{table}

\subsection{Commitment Plans and the Security Level}
\label{app:eval-plans}

Table~\ref{tab:eval-plans} commits the full Llama-2-7B under each of the six candidate plans
(Appendix~\ref{app:plans}) and answers the same 2{,}048-token queries in mode~$\sC$. Codewords of at least $2^{18}$
points open \PlanColsEighteen{} columns per query, against \PlanColsSixteen{} at $2^{16}$ points, for a one-time
encoding of \PlanCommitEighteen{}\,s against \PlanCommitSixteen{}\,s. For Llama-2-7B the planning rule picks
\PlanAutoLlama{}; our runs use $n\ge2^{\PlanLogPoints{}}$ for every model. Table~\ref{tab:eval-lambda} varies the security level:
at $\lambda=40$ a 2{,}048-token proof in mode~$\sC$ is \LamFortyProofC{}\% of its size at $\lambda=128$ (a $\Kpre$
proof, the claims alone, does not change), and the CPU verifier takes \LamFortyVerifyC{}\% ($\sC$) and
\LamFortyVerifyK{}\% ($\Kpre$) of its time.

\begin{table}[t]
  \caption{Mode $\sC$'s candidate plans on Llama-2-7B---longer codewords open fewer columns per query, at a higher one-time encoding cost. The same 2{,}048-token queries; encoding on the \PlatformGpu{}, verifier on 8 threads of the \PlatformCpu{}. $n\ge2^e$: every matrix at $\max(2^e,2\,\mathrm{next\_pow2}(k))$ points; tight: the per-operation lengths; rate~$R$: every matrix at rate $R$; all with the transposed layout and lookup tables}
  \label{tab:eval-plans}
  \centering\footnotesize
  \setlength{\tabcolsep}{3pt}
  \begin{tabular}{@{}lrrrrrr@{}}
    \toprule
    Plan & trees & columns & encode (s) & proof (GB) & prove (s) & verify (s) \\
    \midrule
    \pvitable{tab_eval_plans}
    \bottomrule
  \end{tabular}
\end{table}

\begin{table}[t]
  \caption{Llama-2-7B at $\lambda=40$, 80 and 128---a lower security level shrinks mode $\sC$'s opened columns and the verifier's checks, which matter only at short prompts. Prover on one \PlatformGpu{}, verifier on 8 threads of an \PlatformCpu{}}
  \label{tab:eval-lambda}
  \centering\footnotesize
  \setlength{\tabcolsep}{3pt}
  \begin{tabular}{@{}rrrrrrrr@{}}
    \toprule
    & & \multicolumn{3}{c}{$\sC$} & \multicolumn{3}{c}{$\Kpre$} \\
    \cmidrule(lr){3-5}\cmidrule(l){6-8}
    $T$ & $\lambda$ & prove & verify & proof & prove & verify & proof \\
    \midrule
    \pvitable{tab_eval_lambda}
    \bottomrule
  \end{tabular}
\end{table}

\subsection{V1 Against Published Systems}
\label{app:eval-compare-v1}

\begin{table}[t]
  \caption{Proof sizes against published systems---V1 shrinks our proof, which stays far larger than those of zkLLM and zkGPT. Ours: mode $\sC$, full models; v0 sends \CompareVOneOptS{}\,MB and \CompareVOneOptM{}\,MB for the two OPT rows}
  \label{tab:eval-compare-v1}
  \centering\footnotesize
  \setlength{\tabcolsep}{4pt}
  \begin{tabular}{@{}llrrr@{}}
    \toprule
    System & Model ($T$) & theirs & ours (v0) & V1 \\
    \midrule
    \pvitable{tab_eval_compare_v1}
    \bottomrule
  \end{tabular}
\end{table}

\begin{table*}[t]
  \caption{Full comparison of v0 with published systems---our prover is faster in every row. Settings and notes ($^a$--$^d$) as in Table~\ref{tab:eval-compare}; rows ``zkLLM, \PlatformGpu{}'': zkLLM's public demo on our GPU (prover only; Appendix~\ref{ext:disc-options}) against our interactive cells}
  \label{tab:eval-compare-full}
  \centering\footnotesize
  \setlength{\tabcolsep}{2.5pt}
  \begin{tabular}{@{}llrrrrrr@{}}
    \toprule
    & & \multicolumn{2}{c}{prover} & \multicolumn{2}{c}{verifier} & \multicolumn{2}{c}{proof} \\
    \cmidrule(lr){3-4}\cmidrule(lr){5-6}\cmidrule(l){7-8}
    System & Model ($T$) & theirs & ours & theirs & ours & theirs & ours \\
    \midrule
    \pvitable{tab_eval_compare_full}
    \bottomrule
  \end{tabular}
\end{table*}

\para{LeNet-5.} Table~\ref{tab:eval-compare} omits two LeNet-5 rows for space (Table~\ref{tab:eval-compare-full}
has them, with VGG-11 and zkLLM on the \PlatformGpu{}). On
LeNet-5, zkCNN (verifier on one core) proves in 441\,ms and verifies in 5.8\,ms with a 71.3\,kB
proof, and Bionetta~\cite{bionetta} proves in 3.8\,s and verifies in 10.0\,ms with an 880\,B proof;
ours proves in \LenetProve{}\,ms and verifies in \LenetVerify{}\,ms with a \LenetProof{}\,kB proof: \LenetProveVsZkcnn{}$\times$ and \LenetProveVsBionetta{}$\times$
faster provers, \LenetVerifyVsZkcnn{}$\times$ and \LenetVerifyVsBionetta{}$\times$ faster verifiers, and a proof \LenetProofVsZkcnn{}$\times$ smaller than
zkCNN's and \LenetProofVsBionetta{}$\times$ larger than Bionetta's.

\para{Proof size.} The implementation improvements leave the proof sizes unchanged. V1 shrinks
the proof where it is largest (Table~\ref{tab:eval-compare-v1}). Only committing the activations and
proving the operations without weights, as zkLLM does, would stop the proof from growing with the
prompt. zkAgent reports about 0.30\,s per token for GPT-2
generation~\cite{zkagent}; Section~\ref{sec:gen} gives our generation costs.

\subsection{V1: Prover Time per Leaf}
\label{app:eval-gkr-leaf}

V1's prover pays for the GKR, about one leaf per value sent: \GkrNsPerLeaf{}\,ns per leaf for a
$2^{25}$-leaf instance on the \PlatformGpu{}, which accounts for \VOneOptSGkr{}\,s of the \VOneOptSProveVOne{}\,s that V1's
prover takes for OPT-125M at 2{,}048 tokens in $\sC$ (Table~\ref{tab:eval-v1}).
